\begin{lemma}[A nonaffine horizontal slice]
Let $F$ be a finite-dimensional real inner-product space, let $1\le C_0<\infty$
and let $v\in\Aclass_{m,C_0}$ in the sense of \noderef{IsFitClass} be centered at
the origin with constant $C_0$ in the sense of \noderef{IsCenteredAt}. If $v$ is
not affine in the sense of \noderef{IsSpaceTimeAffine}, then some horizontal
slice of $v$ has a nonzero midpoint Jensen defect: there are a time $t$ and
points $y_1,y_2\in F$ with
\[
  v\left(\frac{y_1+y_2}{2},t\right)
    \ne\frac{v(y_1,t)+v(y_2,t)}{2}.
\]
\end{lemma}
\begin{proof}
We argue by contraposition: assuming that every horizontal slice of $v$
satisfies the midpoint identity,
\begin{equation}\label{eq:nse-mid}
  v\Bigl(\frac{y_1+y_2}{2},t\Bigr)=\frac{v(y_1,t)+v(y_2,t)}{2}
  \qquad(t\in\R,\ y_1,y_2\in F),
\end{equation}
we show that $v$ is affine in the sense of \noderef{IsSpaceTimeAffine}, which
contradicts the hypothesis.

\emph{Step 0: standing facts about $v$.}
By \noderef{IsFitClass} the function $v$ is continuous on all of space-time and
is a viscosity solution of $F[v]=0$ there; by \noderef{IsLimitSolutionOn}, whose
definition of that notion is precisely the conjunction of the two one-sided
conditions, $v$ is therefore both a viscosity subsolution and a viscosity
supersolution in the senses of \noderef{IsLimitSubOn} and
\noderef{IsLimitSuperOn}. By
\noderef{UniformSubquadratic}, applied with the centering point $(0,0)$, for
every compact $J\subseteq\R$ and every $\varepsilon>0$ there is $M>0$ with
$|v(y,\sigma)|\le\varepsilon|y|^2$ for $|y|\ge M$, $\sigma\in J$; that is, $v$
has subquadratic spatial growth on compact time intervals of $I=\R$ in the sense
of \noderef{SubquadraticGrowth}. Hence \noderef{QuadEnvPathConvex} applies to
$v$ with $I=\R$ and any $T>0$.

\emph{Step 1: each slice is affine.}
Fix $t\in\R$ and put $g(y)=v(y,t)$, a continuous function on $F$. Let
$h(y)=g(y)-g(0)$. From \eqref{eq:nse-mid} with $y_2=0$ we get
$h(y/2)=h(y)/2$ for every $y$, and then \eqref{eq:nse-mid} itself gives
\[
  \frac{h(y_1+y_2)}{2}=h\Bigl(\frac{y_1+y_2}{2}\Bigr)
   =\frac{h(y_1)+h(y_2)}{2},
\]
so $h$ is additive. An additive map is $\mathbb Q$-homogeneous: additivity gives
$h(ny)=nh(y)$ for $n\in\mathbb N$ by induction, $h(0)=0$, $h(-y)=-h(y)$, and
$h(y/n)=h(y)/n$ by applying the first identity to $y/n$. Since $h$ is
continuous and $\mathbb Q$ is dense in $\R$, $h$ is $\R$-homogeneous, hence
$\R$-linear. By the Riesz representation of a continuous linear functional on
the finite-dimensional inner-product space $F$, there is $p(t)\in F$ with
\begin{equation}\label{eq:nse-slice}
  v(y,t)=\langle p(t),y\rangle+c(t)\qquad(y\in F),\qquad c(t):=v(0,t).
\end{equation}

\emph{Step 2: $p$ and $c$ are continuous.}
$c(t)=v(0,t)$ is continuous in $t$ because $v$ is continuous. Fix an orthonormal
basis $e_1,\dots,e_m$ of $F$. By \eqref{eq:nse-slice},
$\langle p(t),e_i\rangle=v(e_i,t)-v(0,t)$, so
$p(t)=\sum_i\bigl(v(e_i,t)-v(0,t)\bigr)e_i$ is continuous in $t$.

\emph{Step 3: the forward envelope of an affine slice.}
Let $\ell>0$ and $(x,t)$ be arbitrary. By \eqref{eq:nse-slice} and
\noderef{QuadEnvP},
\[
  \Penv_\ell v(x,t)
   =\sup_{y\in F}\Bigl\{\langle p(t+\ell),y\rangle+c(t+\ell)
      -\frac{|y-x|^2}{2\ell}\Bigr\}.
\]
Writing $y=x+\ell p(t+\ell)+w$ and expanding, the expression in braces equals
\[
  \langle p(t+\ell),x\rangle+c(t+\ell)+\frac{\ell|p(t+\ell)|^2}{2}
    -\frac{|w|^2}{2\ell},
\]
which is maximal exactly at $w=0$. Hence
\begin{equation}\label{eq:nse-P}
  \Penv_\ell v(x,t)=\langle p(t+\ell),x\rangle+c(t+\ell)
     +\frac{\ell}{2}\bigl|p(t+\ell)\bigr|^2\qquad(\ell>0),
\end{equation}
and by \noderef{QuadEnvP} the same formula holds at $\ell=0$, where it reads
$\Penv_0v(x,t)=v(x,t)=\langle p(t),x\rangle+c(t)$.

\emph{Step 4: $p$ is affine in $t$.}
Fix $t$ and $T>0$. By \noderef{QuadEnvPathConvex}, part (1), applied with
$d=0$, $x_0=x$ and this $T$ — the admissibility hypothesis being vacuous since
$I=\R$ — the function $\ell\mapsto\Penv_\ell v(x,t)$ is convex on $[0,T]$. Let
$\ell_0,\ell_1\in[0,T]$, $\theta\in[0,1]$ and
$\ell_\theta=(1-\theta)\ell_0+\theta\ell_1$. Putting $A(\ell)=p(t+\ell)$ and
$B(\ell)=c(t+\ell)+\tfrac\ell2|p(t+\ell)|^2$, convexity and \eqref{eq:nse-P}
give, for every $x\in F$,
\[
  \bigl\langle A(\ell_\theta)-(1-\theta)A(\ell_0)-\theta A(\ell_1),\,x\bigr\rangle
   \le(1-\theta)B(\ell_0)+\theta B(\ell_1)-B(\ell_\theta).
\]
The right side does not depend on $x$. If the vector
$\Delta=A(\ell_\theta)-(1-\theta)A(\ell_0)-\theta A(\ell_1)$ were nonzero, then
taking $x=\lambda\Delta$ with $\lambda\to+\infty$ would make the left side tend
to $+\infty$, a contradiction. Hence $\Delta=0$:
\[
  A(\ell_\theta)=(1-\theta)A(\ell_0)+\theta A(\ell_1)
  \qquad(\ell_0,\ell_1\in[0,T],\ \theta\in[0,1]).
\]
Since $T>0$ was arbitrary, this holds for all $\ell_0,\ell_1\ge0$. Taking
$\ell_0=0$, $\ell_1=1$ and $\theta=\ell\in[0,1]$ gives
$A(\ell)=A(0)+\ell\bigl(A(1)-A(0)\bigr)$ for $\ell\in[0,1]$; and for $\ell>1$,
taking $\ell_0=0$, $\ell_1=\ell$ and $\theta=1/\ell$ gives
$A(1)=(1-1/\ell)A(0)+(1/\ell)A(\ell)$, i.e. again
$A(\ell)=A(0)+\ell(A(1)-A(0))$. Hence, as $t\in\R$ was arbitrary,
\begin{equation}\label{eq:nse-pstep}
  p(t+\ell)=p(t)+\ell\bigl(p(t+1)-p(t)\bigr)\qquad(t\in\R,\ \ell\ge0).
\end{equation}
Note that \eqref{eq:nse-pstep} is licensed only for \emph{nonnegative} offsets
$\ell$, because the convexity it came from is convexity on $[0,T]$. Passing from
it to a formula valid on all of $\R$ is therefore a step to be carried out, not a
remark, and we record it as a statement about an arbitrary map, because Step 6
needs it a second time.

\emph{Base-point lemma.} Let $V$ be a real vector space and let $q:\R\to V$
satisfy
\begin{equation}\label{eq:nse-nonneg-offset}
  q(t+\ell)=q(t)+\ell\bigl(q(t+1)-q(t)\bigr)\qquad(t\in\R,\ \ell\ge0).
\end{equation}
Then $q(\tau)=q(0)+\tau\bigl(q(1)-q(0)\bigr)$ for every $\tau\in\R$.

To see this, fix $\tau\in\R$ and choose the single base point
$s=\min\{\tau,0\}-1$. Then $s\le\tau-1$ and $s\le-1$, so the three offsets
\[
  \tau-s\ge1,\qquad 0-s\ge1,\qquad 1-s\ge2
\]
are all nonnegative, and each of the three instantiations below is admissible in
\eqref{eq:nse-nonneg-offset}. Put $m=q(s+1)-q(s)$. Applying
\eqref{eq:nse-nonneg-offset} at $t=s$ with offsets $\tau-s$, $0-s$ and $1-s$ in
turn, and simplifying $s+(\tau-s)=\tau$, $s+(0-s)=0$ and $s+(1-s)=1$,
\[
  q(\tau)=q(s)+(\tau-s)m,\qquad
  q(0)=q(s)+(0-s)m,\qquad
  q(1)=q(s)+(1-s)m .
\]
Subtracting the second display from the third gives
$q(1)-q(0)=\bigl((1-s)-(0-s)\bigr)m=m$: the slope $m$ at the base point
\emph{is} $q(1)-q(0)$, which is what the base point being unknown might
otherwise have cost us. Subtracting the second from the first gives
$q(\tau)-q(0)=\bigl((\tau-s)-(0-s)\bigr)m=\tau m$. Combining the two,
$q(\tau)=q(0)+\tau\bigl(q(1)-q(0)\bigr)$, which proves the lemma. No case split
on the sign of $\tau$ occurs, and the case $\tau\ge0$ is covered along with the
rest: every instantiation stays inside the nonnegative range of offsets that
\eqref{eq:nse-nonneg-offset} licenses. This is the sense in which the affine
interpolation identity holds on every interval, rather than only forward from one
fixed base point.

Applying the base-point lemma with $V=F$ and $q=p$, whose hypothesis
\eqref{eq:nse-nonneg-offset} is exactly \eqref{eq:nse-pstep}, and setting
$p_0=p(0)$ and $p_1=p(1)-p(0)$,
\begin{equation}\label{eq:nse-paffine}
  p(\tau)=p_0+\tau p_1\qquad(\tau\in\R).
\end{equation}

\emph{Step 5: $p$ is constant.}
Let $C_g>0$ be the constant of \noderef{CenteredGrowth} for this $C_0$; since
$v\in\Aclass_{m,C_0}$ is centered at the origin, for $K\ge2$,
$|v(\zeta)|\le C_gKL(K)^2$ on $\mathcal B_K$, where $L(K)=1+\log K$. If
$F=\{0\}$ then $p(\tau)\in F=\{0\}$ for every $\tau$, so $p_1=p(1)-p(0)=0$ and
$p$ is constant outright, which is the assertion of this step; assume therefore
$F\ne\{0\}$, so that $F$ contains a unit vector. Fix a unit vector $e\in F$, let
$K\ge2$ and let $|\tau|\le K$. By \eqref{eq:nse-slice},
$v(Ke,\tau)-v(-Ke,\tau)=2K\langle p(\tau),e\rangle$, and both points lie in
$\mathcal B_K$, so
\[
  \bigl|\langle p(\tau),e\rangle\bigr|
   \le\frac{|v(Ke,\tau)|+|v(-Ke,\tau)|}{2K}\le C_gL(K)^2 .
\]
Taking the supremum over unit vectors $e$ gives $|p(\tau)|\le C_gL(K)^2$ for
$|\tau|\le K$. In particular $|p(K)|\le C_g(1+\log K)^2$ for every $K\ge2$. By
\eqref{eq:nse-paffine}, $|p(K)|\ge K|p_1|-|p_0|$; if $p_1\ne0$ this grows
linearly in $K$, while $C_g(1+\log K)^2$ grows slower than linearly, so
$K|p_1|-|p_0|\le C_g(1+\log K)^2$ fails for large $K$. Hence $p_1=0$ and
$p(\tau)\equiv p_0=:p$.

\emph{Step 6: $c$ is convex and concave, hence affine.}
With $p$ constant, \eqref{eq:nse-P} becomes
\[
  \Penv_\ell v(x,t)=\langle p,x\rangle+c(t+\ell)+\frac{\ell}{2}|p|^2 .
\]
Fix $t$ and $T>0$; convexity of $\ell\mapsto\Penv_\ell v(x,t)$ on $[0,T]$, from
\noderef{QuadEnvPathConvex} part (1) as in Step 4, and the fact that adding the
affine function $\ell\mapsto\langle p,x\rangle+\tfrac\ell2|p|^2$ changes neither
convexity nor its failure, show that $\ell\mapsto c(t+\ell)$ is convex on
$[0,T]$. Since $t$ and $T$ are arbitrary, $c$ is convex on $\R$.

For the reverse, compute the backward envelope: by \eqref{eq:nse-slice} and
\noderef{QuadEnvM}, for $\ell>0$,
\[
  \Menv_\ell v(x,t)
   =\inf_{y}\Bigl\{\langle p,y\rangle+c(t-\ell)+\frac{|y-x|^2}{2\ell}\Bigr\}
   =\langle p,x\rangle+c(t-\ell)-\frac{\ell}{2}|p|^2 ,
\]
the infimum being attained at $y=x-\ell p$, as the substitution
$y=x-\ell p+w$ shows. By \noderef{QuadEnvPathConvex}, part (2), applied with
$I=\R$ and any $T>0$ — its hypotheses hold because $v$ is a viscosity
supersolution with the required continuity and growth by Step 0 — the function
$\ell\mapsto\Menv_\ell v(x,t)$ is concave on $[0,T]$. Removing the affine
function $\ell\mapsto\langle p,x\rangle-\tfrac\ell2|p|^2$ shows that
$\ell\mapsto c(t-\ell)$ is concave on $[0,T]$ for every $t$ and $T$, i.e. $c$ is
concave on $\R$.

A function on $\R$ that is both convex and concave is affine. Convexity and
concavity together give equality in every chord inequality, so
\begin{equation}\label{eq:nse-cchord}
  c\bigl((1-\theta)\tau_0+\theta\tau_1\bigr)
   =(1-\theta)c(\tau_0)+\theta c(\tau_1)
  \qquad(\tau_0,\tau_1\in\R,\ \theta\in[0,1]).
\end{equation}
We check that $c$ satisfies the hypothesis \eqref{eq:nse-nonneg-offset} of the
base-point lemma of Step 4, with $V=\R$. Let $t\in\R$ and $\ell\ge0$. If
$0\le\ell\le1$, then $t+\ell=(1-\ell)t+\ell(t+1)$ is a convex combination of $t$
and $t+1$, so \eqref{eq:nse-cchord} with $\tau_0=t$, $\tau_1=t+1$ and
$\theta=\ell$ gives
$c(t+\ell)=(1-\ell)c(t)+\ell c(t+1)=c(t)+\ell\bigl(c(t+1)-c(t)\bigr)$. If
$\ell>1$, then $1/\ell\in(0,1)$ and $t+1=(1-1/\ell)t+(1/\ell)(t+\ell)$, so
\eqref{eq:nse-cchord} with $\tau_0=t$, $\tau_1=t+\ell$ and $\theta=1/\ell$ gives
$c(t+1)=(1-1/\ell)c(t)+(1/\ell)c(t+\ell)$; multiplying by $\ell$ and rearranging
yields the same identity $c(t+\ell)=c(t)+\ell\bigl(c(t+1)-c(t)\bigr)$. So
\eqref{eq:nse-nonneg-offset} holds for $q=c$, and the base-point lemma gives
\begin{equation}\label{eq:nse-caffine}
  c(\tau)=c(0)+\tau\bigl(c(1)-c(0)\bigr)\qquad(\tau\in\R).
\end{equation}

\emph{Conclusion.}
Writing $\beta=c(1)-c(0)$ and $\gamma=c(0)$, \eqref{eq:nse-slice},
\eqref{eq:nse-caffine} and Step 5 give
\[
  v(x,t)=\langle p,x\rangle+\beta t+\gamma\qquad\bigl((x,t)\text{ arbitrary}\bigr),
\]
so $v$ is affine in the sense of \noderef{IsSpaceTimeAffine}. This contradicts
the hypothesis that $v$ is not affine, and therefore
\eqref{eq:nse-mid} must fail for some $t,y_1,y_2$, which is the assertion.
\end{proof}
